\clearpage
\fancyhead[L]{\small KOMET.}
\SourceLeaf{425}
\begin{center}
\Large KOMET.

\medskip
Göttingische gelehrte Anzeigen. Stück 83. Seite 825..829. 24. Mai 1819.
\end{center}

\par\indent Im 28. Stück dieser Anzeigen [S. 417 dieses Bandes] haben wir über die beiden Kometen vom vergangenen Jahr berichtet und dabei bereits auf die Wahrscheinlichkeit der Identität des einen dieses Kometen mit dem ersten von 1805 hingewiesen. Die späteren Rechnungen von Herrn Professor \textsc{Encke} haben diese Wahrscheinlichkeit so sehr erhöht, dass an der Identität kaum noch ein Zweifel bestehen kann. Wir zählen diese Erweiterung unserer Kenntnis des Sonnensystems ohne Bedenken zu den bemerkenswertesten Entdeckungen des neunzehnten Jahrhunderts und teilen hier das Wesentliche darüber aus einigen Briefen dieses Astronomen an Herrn Hofrat \textsc{Gauss} mit. Da außer den in der genannten Ausgabe abgedruckten Beobachtungen des Kometen von 1818 noch immer keine weiteren brauchbaren bekannt geworden sind, hielt es Herr Professor \textsc{Encke} vorerst für das Ratsamste, die Beobachtungen des Kometen von 1805 erneut zu untersuchen. Dabei zeigt sich das überraschende Ergebnis, dass eine Ellipse mit einer Umlaufzeit von 3.42 Jahren die Beobachtungen am besten darstellt. Für den Kometen von 1818 hatte er aus den angeführten Elementen eine Umlaufzeit von 3.59 Jahren gefunden. Beide Angaben waren natürlich noch mit einer gewissen Unsicherheit behaftet. Da zwischen den beiden beobachteten Zeiten des Durchgangs durch die Sonnennähe 13.19 Jahre vergangen sind, war es unter der Voraussetzung der Identität am wahrscheinlichsten, dass der Komet in der Zwischenzeit vier Umläufe vollzogen hatte (Dass der Komet, der zwar nicht zu den lichtschwächsten gehört, aber nur mit einem Fernrohr sichtbar ist, bei seiner dreimaligen Rückkehr zur Sonnennähe in den Jahren 1809, 1812 und 1815 unbemerkt blieb, kann niemanden befremden, der mit der Sache vertraut ist). Mit einer entsprechend abgeänderten Umlaufzeit berechnete Herr Professor \textsc{Encke} neue elliptische Elemente aus den Beobachtungen von 1805, die sich damit sehr befriedigend darstellen ließen. Es war eine angenehme Überraschung, dass die einzelnen Elemente dieser neuen Bahn denen der aus den Beobachtungen von 1818 berechneten Bahn außerordentlich viel näher kamen. Nach diesem ermutigenden Erfolg entschloss sich Herr \textsc{Encke}, nach der Methode von Herrn Hofrat \textsc{Gauss} die Störungen zu berechnen, die jene Elemente während der ganzen 13 Jahre durch die Einwirkung des Jupiter erfahren haben mussten. Auch diese mühsame Arbeit blieb nicht unbelohnt, denn sie führte zu einer noch besseren Übereinstimmung. Wir geben hier die Elemente so an, wie sie Herr \textsc{Encke} nach Berücksichtigung der Einwirkungen des Jupiter für den Anfang des Jahres 1819 berechnet hat. Sie beruhen allein auf den Beobachtungen von 1805 und wurden nur in Bezug auf die große Achse so bestimmt, dass die Zeit des Durchgangs durch die Sonnennähe nach vier Umläufen mit der 1819 beobachteten zusammenfiel:

\begin{center}
\begin{tabular}{l l}
\multicolumn{2}{l}{Durchgang durch die Sonnennähe 1819 Jan. 27.27} \\
Länge der Sonnennähe . . . . . . . . . . . . & $156^\circ$ $59'$ $30''$ \\
Länge des aufsteigenden Knoten . . . . . . . . . & 334\ \ 31\ \ \ 0 \\
\multicolumn{2}{l}{\quad beide vom mittleren Äquinoktium von 1819 an gezählt} \\
Neigung der Bahn . . . . . . . . . . . . . . & 13\ \ 36\ \ 35 \\
Exzentrizitätswinkel . . . . . . . . . . . . . & 58\ \ \ 2\ \ 58 \\
Logarithmus der großen Halbachse . . . . . . . . & 0.3451979 \\
\end{tabular}
\end{center}

\Needspace{12\baselineskip}\par\indent Schließlich bestimmte Herr \textsc{Encke} die Bahn noch einmal allein aus den Beobachtungen \SourceLeaf{426}von 1818..1819. Dabei setzte er den Logarithmus der großen Halbachse $= 0.345$ voraus und fand, dass sich die Beobachtungen damit noch viel besser darstellen ließen als mit der früher berechneten Bahn. Die neuen Elemente sind folgende:

\begin{center}
\begin{tabular}{l l}
Zeit der Sonnennähe 1819 Januar . . . . . . . . & 27.27545 Seeberger Zeit \\
Länge der Sonnennähe . . . . . . . . . . . . & $157^\circ\,5'\,53$: \\ 
Länge des Knoten . . . . . . . . . . . . . . & $334\;\;43\;\;37$ \\
Neigung der Bahn . . . . . . . . . . . . . . & $13\;\;38\;\;42$ \\
Exzentrizitätswinkel . . . . . . . . . . . . . & $58\;\;6\;\;45.52$ \\
\end{tabular}
\end{center}

\par\indent Diese Elemente stimmen mit den aus den Beobachtungen von 1805 abgeleiteten so gut überein, dass die Wahrscheinlichkeit der Identität beider Kometen dem gleichkommt, was man bei Gegenständen außerhalb der reinen Mathematik Gewissheit nennt.

\par\indent Man könnte nun noch genauer wissen wollen, inwieweit diese herrliche Übereinstimmung von der vorausgesetzten Zahl von \emph{vier} Umläufen abhängt. Es scheint uns wichtig, auch die Bahnen unter der Annahme von nur drei Umläufen zu bestimmen. Es lässt sich voraussehen, dass sich sowohl die Beobachtungen von 1805 als auch die von 1818 durch Ellipsen mit einer Umlaufzeit von 4.4 Jahren noch recht gut darstellen lassen. Nach einer vorläufigen Abschätzung von Herrn Hofrat \textsc{Gauss} wird die Übereinstimmung beider Bahnen aber erheblich geringer sein als bei den oben angeführten. Wir hoffen, dass sich Herr \textsc{Encke}, der sich durch seine bisherigen Rechnungen schon ein so großes Verdienst erworben hat, auch dieser Untersuchung annehmen wird, und wir neigen zu der Annahme. dass das Ergebnis so ausfallen wird, dass sich auch ohne vollständige Berechnung der Störungen über die Unhaltbarkeit der Annahme einer anderen Zahl von Umläufen als vier urteilen lässt. Diese Untersuchung würde sich übrigens erübrigen, wenn man den Kometen in diesem Jahr noch einmal beobachten könnte. In unseren nördlichen Gegenden besteht dazu freilich wenig oder gar keine Hoffnung. In Italien scheint es jedoch bei Verwendung sehr lichtstarker Instrumente nicht ganz unmöglich zu sein. Um solche Versuche zu erleichtern, teilen wir daher hier noch eine von Herrn \textsc{Encke} berechnete Ephemeride für die kommenden Monate mit. Die Zeiten sind $4^{\mathrm{u}}\,26^{\mathrm{m}}$ Minuten Seeberger Zeit.
\nopagebreak[4]
\begin{center}
\begin{tabular}{c|c|c|c}
1819 & Rektaszension & Südliche & Log. der Entf. \\
 & & Deklination & von der Erde \\
\hline
Mai\quad 31 & $336^\circ\,28'$ & $19^\circ\,43'$ & 0.2277 \\
Juni\quad 10 & $335\;\;27$ & $20\;\;22$ & 0.2168 \\
\phantom{Juni\quad} 20 & $333\;\;30$ & $21\;\;19$ & 0.2070 \\
\phantom{Juni\quad} 30 & $330\;\;36$ & $22\;\;28$ & 0.2004 \\
Juli\quad 10 & $326\;\;49$ & $23\;\;43$ & 0.1991 \\
\phantom{Juli\quad} 20 & $322\;\;26$ & $24\;\;54$ & 0.2049 \\
\phantom{Juli\quad} 30 & $317\;\;47$ & $25\;\;51$ & 0.2189 \\
Aug.\quad 9 & $313\;\;17$ & $26\;\;30$ & 0.2408 \\
\phantom{Aug.\quad} 19 & $309\;\;20$ & $26\;\;47$ & 0.2691 \\
\phantom{Aug.\quad} 29 & $306\;\;11$ & $26\;\;48$ & 0.3020 \\
Sept.\quad 8 & $303\;\;54$ & $26\;\;35$ & 0.3374 \\
\phantom{Sept.\quad} 18 & $302\;\;29$ & $26\;\;13$ & 0.3737 \\
\end{tabular}
\end{center}

\Needspace{7\baselineskip}\par\indent Dabei ist noch zu bemerken, dass diese Ephemeride aus den früheren Elementen berechnet wurde. Die folgenden zwei aus den neueren Elementen berechneten Positionen geben eine Vorstellung davon, wie weit diese Elemente von den früheren abweichen:

\begin{center}
\begin{tabular}{c|c|c|c}
Juli\quad 30 & $320^\circ\,32$ & $24^\circ\,36'$ & 0.2097 \\
Aug.\quad\;\;9 & $315\;\;59$ & $25\;\;21$ & 0.2294 \\
\end{tabular}
\end{center}

\SourceLeaf{427}\par\indent Wir schließen diesen Bericht mit dem Wunsch, dass die höchst bemerkenswerte Entdeckung die Freunde der Astronomie anspornen möge, mit doppeltem Eifer nach Kometen zu suchen, denn alles lässt vermuten, dass sie nur der Anfang einer unermesslichen und nach und nach reifenden Erndte sein wird.
